\chapter{技能编译与认知沉降}
\label{chap:skill-compilation}

\begin{chapteroverviewbox}
本章讨论 AgentOS 从“能够完成任务”走向“形成能力”的关键动力学机制：技能编译（Skill Compilation）与认知沉降（Cognitive Sedimentation）。我们将看到，一个成熟智能体不应在每一次面对熟悉问题时都重新占用工作记忆 $h(t)$、重新展开显式推理、重新构造动作序列；相反，经过反复验证的显式认知过程，应当逐步经历
\[
h
\rightarrow
E
\rightarrow
\Theta
\rightarrow
\mathrm{BodySkill}
\]
的结构转化，使原本需要高层认知介入的 Agent-CSO 逐渐迁移到更加局部、快速和稳定的功能闭环中。这个过程并不是简单的“记住答案”，也不是把一段推理缓存起来，而是把预测结构、控制策略、适用条件、残差检测机制和重新认知入口共同编译成一个可执行闭环。

因此，本章将在此前 CSO/Agent-CSO、三相记忆、工作记忆有限性、Body Runtime、BPCC 以及“异常向上、能力向下”的多尺度原则之上，建立技能形成的动力学描述。我们特别强调：所谓“认知下沉”并不是复杂度像物质一样向低层流动，而是同一类认知结构逐渐改变其表示方式、功能承载位置、执行时间尺度和显式程度。其本体描述应当从早期的“复杂度下沉”进一步修正为
\[
\boxed{
\mathrm{Explicit\ Agent\mbox{-}CSO}
\longrightarrow
\mathrm{Procedural\ Agent\mbox{-}CSO}
}
\]
而所谓成熟能力，则是这种迁移在现实反馈下形成的稳定结果。
\end{chapteroverviewbox}

\section{从显式认知控制到可编译能力}

我们首先从一个最常见、但也最容易被忽略的智能现象开始：同一个智能体在第一次完成某项任务和第一千次完成该任务时，其内部认知组织方式不应该相同。第一次进入陌生环境时，智能体必须把大量状态提升进入工作记忆 $h(t)$：它需要理解当前情境、识别目标、比较候选动作、构造预测、调用情景记忆、读取结构记忆，并不断根据现实反馈修正计划。此时，一个动作的产生可能依赖一个相当长的高层闭环：
\[
World
\rightarrow
SGC/FCS
\rightarrow
h
\rightarrow
Prediction
\rightarrow
Planning
\rightarrow
Action
\rightarrow
World'.
\]
如果一个长期运行的 Agent 在任务已经高度熟练之后仍然重复这一过程，那么从 AgentOS 的角度看，这不是“谨慎”，而是一种结构上的不成熟。原因在于工作记忆具有有限容量，显式认知属于昂贵资源。每一个已经稳定掌握、却仍然持续占据 $h(t)$ 的任务结构，都会挤占新的、真正需要全局推理的问题所需要的认知空间。因此，技能形成首先是由工作记忆有限性推导出的必然要求，而不是一个附加的性能优化。

我们可以把显式处理某一认知结构 $C$ 时的 Agent-CSO 写成
\[
\mathcal C^{\mathrm{exp}}_A(t)
=
\left(
C,
R_h,
F_{\mathrm{cog}},
\tau_h,
B_{\mathrm{kernel}},
S_{\mathrm{active}}
\right),
\]
其中 $R_h$ 表示结构目前以适合显式推理的形式存在，$F_{\mathrm{cog}}$ 表示它由中央认知功能承担，$\tau_h$ 表示其主要活动位于工作记忆时间尺度，$B_{\mathrm{kernel}}$ 表示其承载边界位于认知内核，而 $S_{\mathrm{active}}$ 表示该结构必须显式激活后才能完成任务。技能形成并不要求本体对象 $C$ 消失，而是要求这些承载坐标发生变化：
\[
(R_h,F_{\mathrm{cog}},\tau_h,B_{\mathrm{kernel}},S_{\mathrm{active}})
\longrightarrow
(R_p,F_{\mathrm{local}},\tau_p,B_{\mathrm{body}},S_{\mathrm{procedural}}),
\]
通常满足
\[
\tau_p \ll \tau_h.
\]
因此所谓“熟练”，首先是一种 Agent-CSO 的功能位置迁移和时间尺度迁移。智能体并没有失去关于任务的结构，而是降低了该结构每次执行时所要求的显式认知参与。

这也让我们能够严格地区分“会做”与“形成技能”。一次成功只能说明当前高层认知系统能够构造一条成功轨迹：
\[
\Gamma^{(1)}:
x_0\rightarrow x_1\rightarrow\cdots\rightarrow x_T.
\]
形成技能则要求智能体发现，在一类相似状态集合 $\mathcal X_s$ 中，可以存在一个比重新进行完整显式推理更加局部的映射：
\[
\pi_s:
\mathcal X_s
\rightarrow
\mathcal U_s,
\]
并且该映射能够在给定现实约束和允许残差范围内反复闭合。换言之，技能不是“某个动作序列被记住”，而是系统从多次显式求解中抽取出一个可重用的局部闭环。

因此，本章所说的编译并不是计算机程序意义上的简单语法转换，而是一种跨认知尺度的结构编译：
\[
\boxed{
\mathrm{Explicit\ Cognitive\ Control}
\rightarrow
\mathrm{Repeated\ Experience}
\rightarrow
\mathrm{Stable\ Generative\ Structure}
\rightarrow
\mathrm{Executable\ Local\ Loop}
}
\]
它使一个成熟 Agent 的高层认知逐渐从“亲自完成每个动作”转变为“只在需要时建立结构、设定约束并处理异常”。这也是我们此前提出
\[
\boxed{
\mathrm{Mature\ Intelligence}
=
\mathrm{Selective\ Explicit\ Cognition}
}
\]
这一命题的具体工程含义：真正成熟的智能，不是所有时候都在进行最多的显式思考，而是能够准确判断哪些问题已经不再值得占据显式认知。

\section{重复轨迹：经验如何形成可编译结构}

显式认知过程能够下沉的前提并不是“重复次数足够多”，而是重复经验中逐渐显现出具有稳定结构的轨迹族。对于任务 $g$，设 Agent 在不同时间产生一组现实验证后的经验轨迹
\[
\mathcal E_g
=
\left\{
\Gamma_g^{(1)},
\Gamma_g^{(2)},
\ldots,
\Gamma_g^{(N)}
\right\},
\]
其中每条轨迹都可以包含观测、内部状态、控制参考、动作、结果和残差：
\[
\Gamma_g^{(k)}
=
\left[
x_t,
z_t,
u_t,
y_{t+1},
\varepsilon_{t+1}
\right]_{t=t_0}^{t_1}.
\]
这里的关键不在于轨迹完全相同。现实环境不会机械重复，同一个技能在不同初始条件、不同负载、不同对象和不同扰动下通常会呈现不同的具体轨迹。真正需要识别的是这些轨迹之间是否存在一个可复用的低维结构，即是否能够从大量具体经历中抽取一个稳定的控制不变量。

因此，我们不能简单定义
\[
N>N_0
\Rightarrow
\mathrm{Compile}.
\]
更合理的编译候选条件应当依赖至少若干统计量。我们可以定义任务族 $g$ 的重复稳定度
\[
R_g
=
\Phi
\left(
N_g,
\sigma_{\Gamma},
\bar{\varepsilon}_g,
\sigma_{\varepsilon,g},
P_{\mathrm{success}},
M_{\mathrm{safety}}
\right),
\]
其中 $N_g$ 表示有效重复次数，$\sigma_{\Gamma}$ 描述轨迹族在结构意义上的离散程度，$\bar{\varepsilon}_g$ 与 $\sigma_{\varepsilon,g}$ 分别表示长期平均残差和残差方差，$P_{\mathrm{success}}$ 表示成功概率，$M_{\mathrm{safety}}$ 表示距离安全边界的裕度。当
\[
R_g>\theta_{\mathrm{compile}}
\]
并持续一个足够长的观察窗口之后，系统才有理由把这一组经验视为潜在的技能结构。

这一过程与三相记忆具有直接关系。最初，每一次解决任务时，工作记忆 $h$ 中出现的是即时认知状态；任务结束后，其中具有生命周期价值的部分写入 $E$，形成带时间方向的经验轨迹：
\[
h
\longrightarrow
E.
\]
随着多个相似情景不断积累，$E$ 中不再只有“发生过什么”，还逐渐出现跨情景稳定关系，于是可以发生
\[
E
\longrightarrow
\Theta.
\]
这里的 $\Theta$ 不只是一个抽象知识库，它开始形成关于任务结构的生成模型：什么状态通常出现、哪些变化可预测、哪些行动能产生什么结果、什么误差意味着局部闭环失效。我们因此可以把技能形成前的结构学习写成
\[
\Theta_g
=
\mathcal A
\left(
E_g^{(1:N)}
\right),
\]
其中 $\mathcal A$ 是结构抽取或压缩算子，而不是简单记忆复制。

尤其需要注意，$E\rightarrow\Theta$ 的作用不是把所有经验细节固化。恰恰相反，真正可编译的能力依赖于遗忘大量偶然细节，只保留在新情境中仍具有预测价值的结构。设每条经验包含共同部分 $I_g$ 和偶然扰动 $\nu_k$：
\[
E_g^{(k)}
=
I_g+\nu_k,
\]
则结构化学习的目标更接近估计
\[
\hat I_g
=
\arg\min_I
\sum_{k}
D(E_g^{(k)},I)
+
\lambda\Omega(I),
\]
其中 $D$ 衡量结构偏差，$\Omega(I)$ 抑制不必要的表示复杂度。这样形成的 $\hat I_g$ 才具备从“具体经历”进入“稳定能力”的资格。

因此，在我们的理论中，重复并不直接产生技能。重复只是为系统提供一个统计窗口，使其能够判断哪些结构属于现实的稳定规律，哪些只是偶然扰动。只有当经验经过现实验证、跨情境比较和结构压缩之后，才可能发生真正的技能编译。换句话说：
\[
\boxed{
\mathrm{Repetition}
\neq
\mathrm{Skill};
\qquad
\mathrm{Stable\ Structure\ Extracted\ From\ Repetition}
\Rightarrow
\mathrm{SkillCandidate}.
}
\]
这一区分对于长期 Agent 极为关键，否则系统会把一次偶然成功、环境漏洞或错误习惯过早编译成局部能力，从而使暂时性偏差沉降成长期结构。

\section{稳定策略：从轨迹记忆到局部闭环}

当大量经验已经形成稳定结构后，下一步并不是保存一条“最佳轨迹”，而是从轨迹集合中形成能够根据实时状态进行调整的局部策略。原因很直接：真实环境始终存在扰动。若所谓技能只是
\[
u_0,u_1,\ldots,u_T
\]
这样一段开放环动作序列，那么只要环境与训练情境稍有偏差，技能就会迅速失效。因此，技能必须从开放环轨迹转化成闭环策略：
\[
\boxed{
u_t
=
\pi_s(\hat x_t,g_s,c_s)
}
\]
其中 $\hat x_t$ 是局部状态估计，$g_s$ 是当前技能目标，$c_s$ 是约束与上下文。

这里产生了一个重要的理论变化。显式认知阶段，高层 Agent 可能维持一个很大的状态集合：
\[
x_t^{\mathrm{global}}
=
\left(
World,
Self,
Goal,
History,
Constraints,
Alternatives,
Plans,\ldots
\right).
\]
而技能化以后，本地策略不应该继续依赖全部全局状态。它需要寻找一个足以闭合该技能的局部状态：
\[
x_t^{s}
=
\phi_s(x_t^{\mathrm{global}}),
\]
满足
\[
\dim x_t^{s}
\ll
\dim x_t^{\mathrm{global}}.
\]
$\phi_s$ 可以被理解为技能特定的状态压缩器。只有与当前技能执行直接相关的自由度被保留，其余认知变量退出本地高频闭环。正是在这个意义上，技能编译实现了认知沉降：它不是让低层复制整个高层世界模型，而是让低层获得恰好足以稳定完成局部任务的状态表示。

我们可以把技能策略的稳定性要求写成一个局部闭环系统。设真实局部动力学为
\[
x_{t+1}
=
f(x_t,u_t,w_t),
\]
其中 $w_t$ 表示扰动；编译后的策略为
\[
u_t=\pi_s(x_t).
\]
则技能至少要求在适用区域 $\mathcal D_s$ 内存在某种稳定性意义。例如存在 Lyapunov 型函数 $V_s(x)$，使得在正常扰动范围 $\mathcal W_s$ 内
\[
\mathbb E
\left[
V_s(x_{t+1})-V_s(x_t)
\right]
\le
-\alpha_s\|e_t\|^2+\beta_s\|w_t\|^2.
\]
这并不意味着所有技能都必须使用经典 Lyapunov 方法实现，而是指出“技能已形成”必须包含某种局部可恢复性和误差收敛性质，而不能只看单次任务成功率。

因此，所谓 Stable Policy 至少同时具备四种结构特征。第一，它具有状态反馈，而不是单纯回放动作；第二，它只依赖比高层认知更局部的状态集合；第三，它在一定扰动范围内能够自行修正；第四，它知道自己的适用边界。一旦超出这一边界，局部闭环必须停止宣称自己能够继续完成任务，而将问题重新升级。

从 Agent-CSO 角度看，这一变化可以写成
\[
C_g:
F_{\mathrm{reasoning}}
\longrightarrow
F_{\mathrm{policy}},
\]
即同一个任务相关认知结构，其主要功能承载从高层显式推理迁移到局部策略执行。然而这里不能误解为高层结构被完全删除。高层仍然保留：
\[
Index(C_g),
\quad
Meaning(C_g),
\quad
GoalRelation(C_g),
\quad
FailureSemantics(C_g).
\]
真正向下迁移的是“如何在常规条件下完成它”的执行结构，而不是关于该能力全部意义的认知结构。因此，成熟技能实际上形成一种上下层分工：高层知道“为什么、什么时候以及为了什么使用这个技能”，低层知道“在当前身体和环境条件下如何快速把它做出来”。

\section{预测编译：把高层理解变成局部前向模型}

技能能够成为真正稳定的闭环，仅有策略还不够。一个成熟局部能力通常还必须具有关于动作后果的快速预测结构。否则本地控制只能不断进行反应式修正，而不能提前判断即将发生的状态变化。因此，我们把技能编译拆成两个相互关联但不能混同的过程：
\[
\boxed{
\mathrm{Predictive\ Compilation}
}
\]
与
\[
\boxed{
\mathrm{Control\ Compilation}.
}
\]
前者解决“如果这样做，接下来会发生什么”，后者解决“为了使状态向目标移动，现在应该做什么”。

在显式认知阶段，高层预测可能具有丰富的语义结构。Agent 会利用 $h$、$E$ 和 $\Theta$ 推断：
\[
p(X_{t+\Delta}\mid X_t,A_t,G_t,\Theta),
\]
预测窗口可能较长，状态变量可能包括社会意义、任务结果、资源变化和反事实路径。然而，当某项技能逐渐成熟时，没有必要让每一个毫秒级动作都调用如此宽广的预测过程。相反，需要从 $\Theta$ 中抽取与该技能局部闭环相关的动力结构：
\[
\hat x_{t+\delta}
=
\hat f_s
\left(
x_t^s,
u_t
\right),
\]
其中
\[
\delta\ll\Delta.
\]
$\hat f_s$ 就是技能内部的快速前向模型。

预测编译本质上是一次模型范围收缩。高层模型关注的是：
\[
\mathcal X_{\mathrm{global}}
\times
\mathcal U
\rightarrow
\mathcal X_{\mathrm{global}},
\]
而局部技能只需要：
\[
\mathcal X_s
\times
\mathcal U_s
\rightarrow
\mathcal X_s.
\]
因此可以定义一个投影
\[
P_s:
\mathcal X_{\mathrm{global}}
\rightarrow
\mathcal X_s,
\]
以及一个局部动力近似
\[
\hat f_s
\approx
P_s\circ f\circ P_s^{-1},
\]
这里的 $P_s^{-1}$ 并不要求真正可逆，它只表达局部状态是高层状态在当前技能相关自由度上的有效表示。工程上，它可以由 learned dynamics、状态估计器、系统辨识模型、局部神经动力学模型或者物理模型共同实现。

预测编译的价值在于，它使技能具备“提前发现自己即将失败”的能力。定义一步预测残差：
\[
\varepsilon_{t+1}^{pred}
=
x_{t+1}^{obs}
-
\hat f_s(x_t,u_t),
\]
以及滑动统计量：
\[
R_t^{pred}
=
\sum_{k=0}^{K}
\gamma^k
\left\|
\varepsilon_{t-k}^{pred}
\right\|_{W}^{2}.
\]
如果
\[
R_t^{pred}<\theta_s^{normal},
\]
则说明现实仍处于技能所学习的局部动力区域；如果持续超过阈值，则说明实际世界与技能模型之间出现结构失配。这样，技能就不只是一个动作生成器，还具备最基本的“我是否仍然理解当前局部现实”的能力。

这里必须保留我们此前确立的 Reality Authority 原则。预测模型无论训练得多好，都不能取代真实观测：
\[
\boxed{
Prediction
\neq
Observation.
}
\]
因此 $\hat f_s$ 只具有预测权，没有现实定义权。实际状态仍必须通过真实世界重新进入 Body Runtime，再形成新的 SGC/FCS 或局部状态估计。技能只能依据预测提高控制速度，而不能把自己的预测结果直接写成“现实已经如此”。这一边界确保所谓技能编译不会使系统逐渐进入一个只和自己预测互动的封闭世界。

因此，从三相记忆角度看，预测编译可以被理解为：
\[
E
\rightarrow
\Theta
\rightarrow
\hat f_s,
\]
即大量经验首先沉降成长期结构，再从长期结构中提取特定任务的快速局部动力模型。它展示了一个非常重要的智能规律：长期慢结构并不是因为“慢”就远离行动；恰恰相反，慢结构可以被重新编译为高速预测能力，从而以新的表示形式回到快速闭环之中。

\section{控制编译：把显式规划转化为局部策略}

如果说预测编译把“对世界的理解”下沉，那么控制编译则把“如何行动”的经验下沉。在显式阶段，Agent 为获得动作 $u_t$ 可能需要经历完整的规划过程：
\[
u_t
=
\Pi_{\mathrm{plan}}
\left(
h_t,E,\Theta,G_t,\Psi,\mathcal C_t
\right).
\]
这里不仅包含当前状态，还可能涉及长期目标、价值约束、候选方案比较和风险权衡。这种计算适合陌生、复杂、具有高度不确定性的情景，却不适合每一个已经熟练的动作。

技能形成以后，应当构造一个局部控制器：
\[
u_t
=
\pi_s
\left(
x_t^s,
r_t^s,
\mathcal E_s
\right),
\]
其中 $r_t^s$ 是上层给出的局部参考，$\mathcal E_s$ 是允许的控制包络。这里出现了我们在 Body Runtime 理论中反复强调的分层原则：
\[
\boxed{
HigherLevelSetsFeasibleRegion,
\qquad
LowerLevelOptimizesInsideIt.
}
\]
高层不再规定每一个细节动作，而是给出目标、允许区域、风险边界和退出条件；局部技能则在这个范围内以更快频率解决具体控制问题。

因此，控制编译不是简单把高层输出蒸馏成一个神经网络，而是把高层显式决策的一部分因果责任重新分配给局部功能环。我们可以定义编译前后高层介入率
\[
\eta_H
=
\frac{N_{\mathrm{high-level\ intervention}}}
{N_{\mathrm{control\ decisions}}}.
\]
对于一个真正成熟、环境稳定的技能，应当看到
\[
\eta_H^{after}
\ll
\eta_H^{before},
\]
同时任务成功率和安全裕度不能下降到不可接受水平：
\[
P_{\mathrm{success}}^{after}
\ge
P_{\mathrm{success}}^{min},
\]
\[
M_{\mathrm{safety}}^{after}
\ge
M_{\mathrm{safety}}^{min}.
\]
这比单纯比较推理 token 数更加准确地描述了技能成熟程度，因为核心变化并不是“算得少了”，而是高层控制责任真正被迁移到了局部闭环。

进一步，我们可以把技能控制器理解为一个受约束优化问题：
\[
u_t^{*}
=
\arg\min_{u}
\left[
L_s
\left(
\hat x_{t+1:t+H},
r_t
\right)
+
\lambda_u C_u(u)
+
\lambda_r R(u)
\right],
\]
满足
\[
x_{t+k}\in\mathcal X_s^{safe},
\qquad
u_{t+k}\in\mathcal U_s^{allowed}.
\]
这与模型预测控制、局部 policy 和传统 servo 的具体实现都可以兼容。AgentOS 的理论并不规定某一种算法，而是规定功能要求：局部技能必须在特定状态空间中利用快速预测和控制，在上层给定的语义目标与约束范围内完成闭环。

尤其重要的是，控制编译必须保留“退出语义”。一个技能不仅需要知道如何继续，还要知道什么时候自己不应该继续。因而实际技能策略不能只有
\[
\pi_s:x\rightarrow u,
\]
还需要一个有效性判别函数
\[
\chi_s(x,\varepsilon,c)
\in
\{0,1\},
\]
其中
\[
\chi_s=1
\]
表示当前仍处于技能可负责的局部区域，
\[
\chi_s=0
\]
则意味着必须中断局部自主控制并把相关结构提升回更高认知层。于是技能本身形成了一个具有权限边界的局部控制主体，而不是无限制接管行动的不可解释黑箱。

从这一点出发我们可以看到，“能力向下”与“异常向上”实际上是同一机制的两个方向。控制编译使稳定的动作责任向下迁移；有效性检测和残差机制则保证一旦环境超出技能经验边界，控制责任可以重新向上返回。这种双向可逆性，正是技能下沉不会破坏长期通用性的关键。

\section{BodySkill：技能作为可执行的局部预测控制闭环}

在完成预测编译和控制编译以后，我们终于可以严格定义 BodySkill。此前我们曾经用一个极简公式表达：
\[
\boxed{
\mathrm{BodySkill}
=
\mathrm{LocalPredictor}
+
\mathrm{LocalController}
+
\mathrm{ResidualModel}.
}
\]
这个公式值得保留，因为它抓住了技能区别于普通记忆、脚本和动作模板的核心结构。Local Predictor 负责快速预测局部状态变化；Local Controller 负责根据目标与预测产生具体动作；Residual Model 负责判断现实是否仍处于技能可解释、可控制的区域。

为了进行更严格的工程形式化，我们可以把技能 $S$ 表示为
\[
S
=
\left(
\mathcal X_S,
\mathcal U_S,
\hat f_S,
\pi_S,
R_S,
\chi_S,
\mathcal C_S,
\mathcal I_S
\right),
\]
其中 $\mathcal X_S$ 是技能局部状态空间，$\mathcal U_S$ 是动作空间，$\hat f_S$ 是局部预测器，$\pi_S$ 是控制策略，$R_S$ 是残差模型，$\chi_S$ 是有效性判别器，$\mathcal C_S$ 是安全和能力约束，而 $\mathcal I_S$ 是高层语义索引，描述这个技能“是什么、什么时候可调用、完成什么意义上的目标”。最后这一项尤其重要，因为如果缺少 $\mathcal I_S$，低层虽然可能具备某种动作能力，高层却无法在适当语义环境中可靠调用它。

因此，BodySkill 不是完整 Agent 的缩小版。它通常不具有完整的 $\Psi$、长期目标系统和开放式生命周期。它只需要形成一个局部闭环：
\[
r_t
\rightarrow
\pi_S
\rightarrow
u_t
\rightarrow
Reality
\rightarrow
x_{t+1}^{obs}
\rightarrow
R_S
\rightarrow
\pi_S.
\]
如果
\[
R_S<\theta_S
\]
且
\[
x_t\in\mathcal D_S,
\]
局部环可以持续自治；一旦违反条件，就触发：
\[
\mathrm{Escalate}(S,\varepsilon,\mathrm{context}).
\]
由此我们避免了一个危险误区：并不是下沉得越深越好，也不是自主性越强越成熟。真正成熟的局部技能具有清晰的能力域、约束域和退出域。

在数字 Body 中，BodySkill 并不要求对应机械动作。例如 Browser 操作、文件处理、Shell 工作流、GUI 操作、API 交互同样可以形成局部技能。假设一个 Agent 初次完成某软件操作需要十余轮视觉观察、语义推理和动作规划，随着重复经验积累，它可以逐步形成
\[
\mathrm{VisualReasoning}
\rightarrow
\mathrm{StableUIStateModel}
\rightarrow
\mathrm{CompiledInteractionSkill}.
\]
以后高层只需给出“导出当前报告”这一语义目标，本地技能即可完成菜单定位、文件路径选择、状态确认等低层步骤。只有当软件界面改变、按钮消失、权限异常或者返回结果与预测显著不一致时，问题才重新进入显式认知。

同理，在物理机器人中，抓取、行走、开门、使用工具等能力也可以由同一种理论描述。具体实现当然会使用不同的状态空间和控制器，但其功能结构仍然保持：
\[
\boxed{
Prediction
\rightarrow
Control
\rightarrow
Reality
\rightarrow
Residual
\rightarrow
Correction.
}
\]
这再次体现了我们在认识论部分提出的动力学同构原则：数字身体与物理身体可以拥有不同材料和控制频率，却形成同类的技能闭环。因此，BodySkill 既是认知沉降的工程终点之一，也是实现 Body Scale 的关键单位。

\section{技能形成的判据：什么时候才允许认知下沉}

如果一个 Agent 只要重复几次行为就自动将其编译成技能，那么长期运行后必然会积累大量错误习惯。因此，技能编译必须被理解成一种受治理的慢变量改变，而不是对短期成功的即时响应。我们需要建立一个明确的编译判据，使系统能够区分“暂时可重复”和“已经足以成为稳定能力”。

设候选技能 $S$ 的评估向量为
\[
\mathbf q_S
=
\left[
q_{\mathrm{rep}},
q_{\mathrm{pred}},
q_{\mathrm{ctrl}},
q_{\mathrm{safe}},
q_{\mathrm{gen}},
q_{\mathrm{recover}}
\right],
\]
分别对应重复性、预测可靠性、控制可靠性、安全裕度、跨情境泛化能力和失败后的可恢复性。可以进一步定义综合可信度
\[
Q_S
=
\sum_i
w_i q_i
-
\lambda C_{\mathrm{local}}
-
\mu R_{\mathrm{systemic}},
\]
其中 $C_{\mathrm{local}}$ 表示编译后维护该技能的结构成本，$R_{\mathrm{systemic}}$ 表示该技能可能给整个 Agent 带来的系统性风险。只有在
\[
Q_S>\theta_{\mathrm{on}}
\]
持续足够长时间之后，系统才允许技能进入正式下沉过程。

这里最好使用迟滞而不是单阈值。设技能激活阈值为 $\theta_{\mathrm{on}}$，保留阈值为 $\theta_{\mathrm{off}}$：
\[
\theta_{\mathrm{on}}
>
\theta_{\mathrm{off}}.
\]
新技能必须拥有较强证据才能形成，但一个已经成熟的技能不会因为短时间性能波动就立即被删除。这与我们此前对慢变量保护和拓扑迟滞的讨论完全一致。技能作为一种相对稳定结构，其变化速度原则上应慢于一次具体任务状态：
\[
\tau_{\mathrm{task}}
\ll
\tau_{\mathrm{skill}}.
\]
否则 Agent 会不停地在“显式—自动—显式—自动”之间抖动，形成认知控制层面的结构震荡。

不过，迟滞并不意味着技能僵化。技能仍然需要持续接收来自现实世界的统计反馈。定义长期性能估计
\[
\bar J_S(t)
=
(1-\alpha)\bar J_S(t-1)
+
\alpha J_S(t),
\]
以及长期残差
\[
\bar R_S(t)
=
(1-\beta)\bar R_S(t-1)
+
\beta R_S(t).
\]
如果技能长期出现
\[
\bar J_S\downarrow,
\qquad
\bar R_S\uparrow,
\]
则说明原先的局部动力结构可能已经不再适用。这时并不意味着立即删除技能，而是首先降低其自治范围、提高高层监督比例，并重新收集显式经验。这符合我们提出的 Minimum Necessary Reorganization 原则：先通过较小代价恢复可控性，再决定是否需要重新学习甚至废弃该能力。

编译判据还必须与风险级别相关。同样的成功概率，对于排版文档可能已经足够，但对于高速车辆制动显然远远不够。因此阈值本身必须是上下文函数：
\[
\theta_{\mathrm{compile}}
=
\Theta
\left(
Risk,
Irreversibility,
SafetyMargin,
Observability
\right).
\]
风险越高、动作越不可逆、现实状态越难充分观测，下沉所需要的验证强度就越高。这一点防止我们把“认知沉降”误解成一种无条件追求速度的优化策略。AgentOS 所追求的不是最大程度自动化，而是**在可验证、可恢复和受约束的前提下，把已经不值得高层反复处理的问题迁移到更合适的时间尺度。**

因此，技能成熟应当定义为一组结构条件同时成立，而不是一个单一成功率指标：
\[
\boxed{
\mathrm{SkillMaturity}
=
\mathrm{Repeatability}
+
\mathrm{Predictability}
+
\mathrm{Controllability}
+
\mathrm{Safety}
+
\mathrm{Generalization}
+
\mathrm{Recoverability}.
}
\]
只有满足这些条件，认知沉降才真正降低了系统未来的认知成本，而不是把今天的错误隐藏进明天的自动化结构。

\section{失败后的重新认知：技能编译必须保持可逆}

技能编译最重要的性质之一是可逆性。一个成熟智能系统不会因为某项能力已经自动化，就永久失去对该能力重新进行显式推理的能力。现实世界始终可能产生新物体、新规则、新故障、新身体状态或者新的目标约束，使原来的技能结构失效。因此
\[
\mathrm{Explicit}
\rightarrow
\mathrm{Procedural}
\]
不能是一条单向不可逆的链，而必须存在反向路径：
\[
\boxed{
\mathrm{Procedural}
\leftrightarrow
\mathrm{Explicit}.
}
\]

设技能 $S$ 正在局部执行，其残差序列为
\[
\varepsilon_{1:t}^{S}.
\]
单次微小残差通常可以由局部控制器自行吸收，因此不应立即触发高层思考，否则系统会重新陷入所有细节都占据 $h$ 的低效状态。我们可以定义升级强度
\[
\Gamma_t^S
=
F
\left(
\|\varepsilon_t\|,
Persistence_t,
Novelty_t,
Risk_t,
GoalRelevance_t,
ControlMargin_t
\right).
\]
当
\[
\Gamma_t^S<\theta_{\mathrm{esc}}
\]
时：
\[
\mathrm{LocalCorrection}.
\]
只有当
\[
\Gamma_t^S\geq\theta_{\mathrm{esc}}
\]
时才发生
\[
\boxed{
\mathrm{ProceduralAgentCSO}
\rightarrow
\mathrm{ExplicitAgentCSO}.
}
\]

重新显式化并不意味着把底层所有状态原样推送到 $h$。这会破坏工作记忆有限容量原则。合理做法是让 Body Runtime 或技能层首先形成异常摘要：
\[
\fitdisplay{
\zeta_S
=
Compress
\left(
Goal,
ExpectedState,
ObservedState,
Residual,
RecentActions,
Risk,
Confidence
\right),
}
\]
然后：
\[
\zeta_S
\rightarrow h.
\]
高层得到的不是毫秒级控制历史，而是一个具有语义意义的认知事件，例如：“抓取策略连续三次预测滑移小于实际滑移，夹持稳定裕度下降，当前局部模型可信度不足。”只有在需要进一步诊断时，高层才向下请求更多局部细节。这样才能实现我们此前概括的：
\[
\boxed{
RealityDoesNotContinuouslyBecomeThought;
\quad
RelevantDeviationBecomesThought.
}
\]

一旦问题重新进入 $h$，Agent 可以利用 $E$ 召回过去类似异常，通过 $\Theta$ 重新构造因果解释，形成新的局部策略，并产生新的现实实验。于是失败不只是技能中断，而成为新的经验生成源：
\[
SkillFailure
\rightarrow
ExplicitReCognition
\rightarrow
NewExperience
\rightarrow
\Theta'
\rightarrow
Skill'.
\]
这意味着“自动化”本身仍然处于生命周期学习之中。局部技能并没有被永久冻结，而是通过异常事件保持与高层认知连接。

这种机制还解决了一个长期持续学习中的核心矛盾：如果所有能力永远停留在显式层，工作记忆会随着能力增长越来越拥挤；如果能力下沉后不能重新显式化，系统又会越来越僵化。真正的长期智能必须同时具备
\begin{statementbox}
DownwardCompilation
\end{statementbox}
和
\begin{statementbox}
UpwardReactivation.
\end{statementbox}
前者提供效率，后者提供适应性。两者之间由现实残差控制，而不是由固定时间表控制。

因此，熟练并不等于“不再思考”，更准确地说是：
\begin{statementbox}
只有现实重新证明值得思考时，才重新思考。
\end{statementbox}
这句话揭示了技能编译与工作记忆有限理论之间的深层关系。成熟 Agent 不是逐渐停止认知，而是逐渐学会把显式认知保留给真正尚未闭合的问题。

\section{从“复杂度下沉”到 Agent-CSO 的程序化迁移}

在早期理论中，我们曾将技能形成描述成“复杂度向低层下沉”。这一表述具有非常好的直觉意义：一个初学任务原本需要大量注意、推理、搜索和计划，熟练以后高层认知负担降低，看起来像是原有复杂度被迁移到了某个更低、更快的层级。这一语言仍然可以作为工程描述保留，但经过 CSO 与 Agent-CSO 本体论的发展，我们需要对它进行更严格的修正。

真正发生的不是某种独立“复杂度实体”从 $h$ 流向 Body。更准确地说，是认知结构的承载坐标发生变化。设某一任务结构 $C_s$ 在学习早期表示为
\[
AgentCSO^{(0)}_A(C_s)
=
\left(
C_s,
R_{\mathrm{explicit}},
F_h,
\tau_h,
B_{\mathrm{kernel}},
S_{\mathrm{active}}
\right).
\]
经过经验积累、结构抽象和编译以后：
\[
AgentCSO^{(1)}_A(C_s)
=
\left(
C_s,
R_{\mathrm{procedural}},
F_{\mathrm{skill}},
\tau_s,
B_{\mathrm{body}},
S_{\mathrm{latent/automatic}}
\right),
\]
其中通常有
\[
\tau_s\ll\tau_h.
\]
如果我们定义迁移向量
\[
\mathcal M_{C_s}
=
\left(
M_R,
M_F,
M_\tau,
M_B
\right),
\]
那么技能学习正是同时发生了表示迁移、功能迁移、时间尺度迁移和承载边界迁移。

这种表述还解释了为什么
\[
\boxed{
ThinkingLess
\neq
KnowingLess.
}
\]
当一个专家不再需要逐步思考一个熟练动作时，并不意味着相关认知结构消失。相反，它往往意味着这些结构已经以更加稳定、更快速、更局部的方式被重新承载。因此我们应当用“结构连续性原则”替代早期容易被误解成物理守恒的“复杂度守恒”直觉：
\begin{proposition*}
有效认知结构在成熟过程中通常不是简单消失，而是改变表示与承载位置。
\end{proposition*}
这里仍然允许遗忘、抽象、合并和结构损失，因此它不是严格守恒定律。

从这一角度，技能编译实际上就是 Agent-CSO 的程序化迁移。所谓“程序化”并不限定为传统软件程序，它泛指认知结构从开放式显式推理表示转化成具有稳定输入、局部状态、预测、控制和退出条件的可执行闭环。因此：
\[
\boxed{
\mathrm{Proceduralization}
=
\mathrm{CSORealizationChange}
}
\]
而不是简单
\[
\mathrm{Memory}
\rightarrow
\mathrm{Action}.
\]

这一步对于 AgentOS 理论具有更一般的意义。三相记忆告诉我们结构可以从即时状态沉降到经验与生成结构；BodySkill 则告诉我们，生成结构还可以重新以高速可执行形式进入现实闭环。于是完整的学习链不再只有
\[
h
\rightarrow
E
\rightarrow
\Theta,
\]
而是扩展为
\[
\boxed{
h
\rightarrow
E
\rightarrow
\Theta
\rightarrow
\mathrm{ProceduralCarrier}
\rightarrow
Reality.
}
\]
与此同时，当现实残差持续出现时又存在
\[
Reality
\rightarrow
Residual
\rightarrow
h.
\]
两者共同构成一种跨尺度循环，而不是单向沉降。

因此，本章最终需要把“技能”重新放回智能动力学的总体框架中理解。技能不是一个外围工具，也不是某种预设模块，而是历史显式认知被现实反复验证之后形成的结构沉积。每一个成熟技能，都可以被看成过去大量认知计算在系统功能结构中的一种长期残留。它减少未来任务所需要的显式协调成本，并因此为新的认知活动释放工作记忆。

\section{技能编译作为智能成熟的基本动力学}

现在我们可以把本章的全部讨论收束成一个更一般的成熟理论。一个长期 Agent 的能力增长，不应简单表现为：
\[
Knowledge(t)\uparrow,
\]
因为如果每增加一种知识和能力都要求更多显式认知，那么系统最终会受到工作记忆上界的严格限制。真正可持续的发展必须满足一个更重要的条件：过去已经解决的问题能够不断改变未来解决问题所需的结构成本。

设任务 $g$ 在第 $n$ 次执行时需要的显式认知资源为
\[
C_h^{(n)}(g).
\]
如果系统真正形成技能，则在适用分布保持基本稳定时，我们期待
\[
\frac{\partial C_h^{(n)}(g)}
{\partial n}
<0,
\]
并在成熟之后趋于一个较低的监督成本
\[
C_h^{(n)}(g)
\rightarrow
C_{\mathrm{supervise}}(g),
\]
而不是严格趋向零。这个剩余项非常重要，因为高层仍需要保持语义目标、权限、异常处理和现实验证。因此，成熟能力不是完全脱离认知，而是从持续执行责任转变成低频监督责任。

与此同时，释放出的工作记忆容量可以处理新的、更开放的问题。若总认知容量上界为 $C_{\max}$，已有成熟任务占用显式资源
\[
C_{\mathrm{mature}}
\]
下降，则新的认知空间
\[
C_{\mathrm{novel}}
\leq
C_{\max}-C_{\mathrm{mature}}
\]
相应扩大。由此产生一个极为重要的发育循环：
\[
\boxed{
NovelProblem
\rightarrow
ExplicitCognition
\rightarrow
Experience
\rightarrow
Structure
\rightarrow
Skill
\rightarrow
ReleasedCognitiveCapacity
\rightarrow
HarderNovelProblem.
}
\]
这就是为什么技能形成不是性能工程的附属现象，而是有限认知系统能够持续扩展能力边界的核心机制。

我们甚至可以定义一个粗略的认知成熟指标。设某一时间窗口内系统成功完成的任务价值为 $V_{\mathrm{task}}$，显式认知占用为 $C_h$，高层异常重新介入率为 $\eta_{\mathrm{esc}}$，则可以考察
\[
M_A
=
\frac{V_{\mathrm{task}}}
{C_h+\lambda\eta_{\mathrm{esc}}}
\]
的长期变化。成熟并不意味着简单最大化该式，因为极低的异常升级率也可能来自系统没有正确检测失败，但它揭示了一个方向：在维持现实正确性、安全性和可恢复性的条件下，一个更成熟的 Agent 应该用更少的中央显式认知完成更多已经掌握的任务，把 $h(t)$ 留给真正的新颖性与不确定性。

因此，本章最终得到四条可以作为 AgentOS 技能动力学基本规律的结论。

第一，
\[
\boxed{
RepeatedExperience
\rightarrow
StableStructure.
}
\]
重复经验只有经过现实验证和结构压缩才形成可编译内容。

第二，
\[
\boxed{
StableStructure
\rightarrow
LocalExecutableLoop.
}
\]
稳定认知结构趋向更加局部、更快速和更低显式成本的功能承载。

第三，
\[
\boxed{
UnexpectedResidual
\rightarrow
ExplicitReactivation.
}
\]
局部结构无法解释现实的时候，相关 Agent-CSO 必须重新提升进入显式认知。

第四，
\begin{statementbox}
CompilationMustRemainReversible.
\end{statementbox}
真正的技能不是冻结，而是具有稳定常规闭环和异常重新学习入口的结构。

于是我们可以给出本章的最终定义：

\begin{definition*}
技能编译，是智能体将经过重复现实验证的显式认知结构，
逐步转换为具有局部预测、局部控制和残差检测能力的可执行闭环，
从而降低其未来显式认知占用，同时保留异常重新认知能力的多尺度结构迁移过程。
\end{definition*}

相应地：

\begin{statementbox}
认知沉降，是 Agent-CSO 从快速、昂贵、全局显式的承载方式，
逐渐迁向稳定、局部、低成本和程序化承载方式的生命周期动力学。
\end{statementbox}

因此，从智能动力学的角度看，真正的“熟练”并不是一个 Agent 能够更快地重复曾经的思考，而是它已经不再需要重复那种思考。过去的显式认知已经转化成了现在的结构，而现在的结构又改变了未来认知必须承担什么。这正是我们在整个理论中反复遇到的基本关系：

\[
\boxed{
PastCognition
\rightarrow
PresentStructure
\rightarrow
FutureCognitiveFreedom.
}
\]

一个持续智能体由此才真正具有发育意义：它不仅累积过去，而且不断把过去编译成未来可以依赖的能力；它不是把越来越多经验堆积在工作记忆周围，而是持续改变“下一次还需要思考什么”。这也构成技能编译与认知沉降在 AgentOS 中最根本的理论意义。
